\section{Die Umdefinition des Randwertproblem}
Das Randwertproblem ist eines der größten Probleme des Algorithmus. Es entsthet dadurch, dass Werte, die am Rand des Kartenausschnittes liegen nicht sicher als Berge klassifiziert werden können. Das liegt daran, dass ein in den Kartenbereich hereinragender Hang auch alks Berg erkannt werden würde. Sein höhster Punkt wäre dadruch ein Gipfel, obwohl er nur ein unbedeutender Punkt auf dem Hang ist. Dieses Problem soll dadurch gelöst werden, dass ein Randbereich definiert wird, der in der Hochpunktfindung nicht betrachet wird. \todo{vergleich und ref zu tehorie Randwert} 
\todo{neueinführung min\_prom\/min\_dom}
\todo{lösung\/Pojektrefernezfehler}
\todo{Border width optimieren (immer optimaler wert wäre der Dominanzthreashold)}
\todo{nicht von min prom inklusion von gipfel abhängig machen}
Das Randwert problem als min\_prominence < x und min\_dominanz < y
 objektreverenzfehler